The main motivation behind the work presented in this chapter was to classify CG traffic in a faster and more reliable way, enabling latency-sensitive network processing such as L4S. Initially, we designed and evaluated an unsupervised model based on USAD model,  which demonstrated superior performance compared to a supervised approach, particularly in handling previously unseen games and entirely new CG platforms.

Building on this, we introduced a hybrid architecture that integrates a P4 hardware implementation\footnote{\url{https://github.com/mosaico-anr/P4_NFV_CG_Detector}} on a Tofino chipset for feature extraction directly within the data plane, while keeping the machine learning model execution in the control plane as VNFs. Implementing this architecture posed significant challenges due to the strict limitations inherent in real P4 hardware. Nevertheless, we proposed several adaptations to circumvent these constraints. These adaptations not only enabled hardware execution but also maintained a classification accuracy comparable to that of the software-based implementation, providing valuable insights and feedback for the networking community. % The resulting feature extraction is performed two orders of magnitude faster

%Our future work will integrate our classifier with a L4S architecture and evaluate the gain for CG platforms to respect their latency requirements under challenging network conditions. Also, we will evaluate our architecture with real traffic from an ISP to further assess the performance and classification results in real conditions.